\subsection{The base detector}
The detector we extend scores anomalies in SCADA-based industrial systems by self-supervised online learning~\citep{Wadinger2023}. For each signal it evaluates the conditional cumulative distribution of the current value given the remaining signals under a multivariate Gaussian model, flags the value when this probability lies outside $[\alpha, 1-\alpha]$ with $\alpha \approx 0.00133$, and updates the conditional model online while remaining aware of change points so that legitimate regime shifts are absorbed rather than alarmed. The detector is interpretable: it identifies the responsible signal, the direction of the excursion, and exposes per-signal operating limits. It localizes anomalies but does not classify the failure mode, which is the gap this work addresses. We reuse its conditional-residual stream as the substrate for typing rather than introducing a separate diagnostic pipeline.

\subsection{Sensor-fault taxonomies}
The categories we attribute---bias, drift, accuracy loss (noise), and freezing (constant)---are standard in the sensor-fault literature rather than novel to this work. \citet{Sharma2010} characterize the recurrent fault signatures of deployed sensors, including offset (bias), gain, and stuck-at faults, and study detection methods against them. \citet{Ni2009} likewise organize sensor data faults into a taxonomy spanning calibration offset, drift, and constant (frozen) readings. We adopt this established taxonomy unchanged; our contribution is its streaming, conditional-residual realization, not the categories themselves.

\subsection{Fault-injection protocols}
Because no public stream simultaneously labels all four fault types per signal and per timestep, controlled injection is the established route to per-type evaluation. \citet{Lai2021} introduce a benchmark methodology that injects parameterized anomaly patterns into time series to evaluate detectors under known ground truth. We follow this precedent, injecting the four typed faults into a multivariate process at controlled magnitudes and onsets so that recall and confusion can be measured per type, and we triangulate the synthetic findings against the real-fault deployment described in Section~\ref{Experimental Setup}.

\subsection{Offline fault classifiers}
A large body of work classifies sensor and process faults offline, typically by training supervised or windowed models on labelled fault episodes and extracting features over batches of history. Such methods can exploit rich temporal context but require stored windows, labelled fault data, and a training phase, which conflicts with the streaming, self-supervised, and label-free setting of the base detector. Our classifier instead derives the type from the running conditional residuals already maintained by the detector, in O(1) memory and without history buffers, so that typing is available at the same latency and operational footprint as detection.
